\section{Related Work}

\xhdr{General purpose heuristic design} There has been significant research in designing general-purpose heuristics for speeding up satisficing planning. The first set of approaches are based on simplifying the search problem for example using landmark heuristics \cite{bonet2001planning, richter2010lama}. The next set of approaches aim to include novelty-based exploration in greedy best-first search \cite{xie2012planning, xie2014adding, lipovetzky2017best, teichteil2021boundary}. The latter set of approaches showed state-of-the-art performance (best-first width search \cite{lipovetzky2017best, teichteil2021boundary}, BFWS) in numerous settings. We show that in domains where data is available, it can be more effective to incorporate a learned heuristic into a greedy best-first search procedure.


\xhdr{Learning heuristic search} There have been numerous previous works that attempt to learn search heuristics: Arfaee \etal \cite{arfaee2011learning} propose to improve heuristics iteratively, Virseda \etal  \cite{us2013learning} learn to combine heuristics to estimate graph node distances, Wilt \etal \cite{wilt2015building} and Garrett \etal \cite{garrett2016learning} propose to learn node rankings, Thayer \etal \cite{thayer2011learning} suggest to infer heuristics during a search, and Kim \etal \cite{KimA20} train a neural network to predict graph node distances. These methods generally do not consider the non-i.i.d. nature of heuristic search. Further, Bhardwaj \etal \cite{bhardwaj2017learning} propose \sail, where heuristic learning is framed as an imitation learning problem with cost-to-go oracles. The \sail heuristic uses hand-designed features tailored for obstacle avoidance, with a linear time-complexity in the number of explored grid nodes found to be colliding with an obstacle. Feature-engineering becomes more difficult as we attempt to learn heuristics on diverse graphs such as ones seen in Section \ref{sec:experiments-diversity}, where we may need expert knowledge. Further, heuristics that do not have a constant time complexity in the size of the graph \cite{bhardwaj2017learning, ariki2019fully, terasawa20203d, yonetani2020path, archetti2021neural} generally scale poorly with graph size and hence have constrained use cases. Recent approaches to learning heuristics include Retro* \cite{chen2020retro} by Chen \etal, where a heuristic is learned in the context of AND-OR search trees for chemical retrosynthetic planning. Our work focuses on a more general graph setting.

There has been significant progress on learning heuristics for NP-hard combinatorial optimization problems \cite{NIPS2017_d9896106, NEURIPS2018_8d3bba74, NEURIPS2020_49f85a9e}. Focusing on solving NP-hard problems allows these approaches to design algorithms that are often non-exact and have a relatively large computational budget. This is not the case for methods that focus on polynomial time search, where learning-based methods are bounded by the determinism and time complexity of classical algorithms such as greedy best-first search.



%Still, these heuristic learning methods, due to their time complexities, are impractical for the application in polynomial-time search problems, on which this work focuses.

%\textbf{\phil generally runs in $\mathcal{O}(1)$ per prediction without requiring any hand-engineering. We provide a comprehensive comparison with \sail in Section \ref{sec:experiments-grids}.}

%recent works by Ariki \etal \cite{ariki2019fully} and Teresawa \etal \cite{terasawa20203d} perform convolutions over full-graphs to infer their heuristic values, which  


\xhdr{Learning general purpose search} 
Learning general search policies is a very well-studied research area with a rich set of developments and applications. These include Monte Carlo Tree Search methods \cite{silver2010monte, guez2018learning}, implicit planning methods \cite{deac2020xlvin, karkus2017qmdp, tamar2016value}, and imagination-based planning approaches \cite{pascanu2017learning, weber2017imagination}. Learning search heuristics can be seen as a special case of \textit{general purpose search}, where the search problem is treated as a partially observable Markov decision process with restricted action evaluation (see Section \ref{sec:approach}), and with models running in $\mathcal{O}(1)$ to remain competitive time-complexity-wise on problems where best-first search performs well. \textit{General purpose search} methods do not take into account the above-mentioned constraints, which motivates the development of tailored approaches for learning heuristics \cite{ bhardwaj2017learning, chen2020retro}.

%\jure{I am not sure what precisely are we trying to say here and how does our method relate/differ from this} \update{The goal here was to answer the question of 'why not just use MCTS or something implicit like VIN'. I tried to rephrase the paragraph so that it is clear. I am also happy to drop this in case you think it is not necessary.}

\xhdr{Imitation learning} Our approach builds on prior work in imitation learning (IL) with cost-to-go oracles. Cost-to-go oracles have been incorporated in the context of IL in methods such as SEARN \cite{daume2009search}, AggreVaTe \cite{ross2014reinforcement}, LOLS \cite{chang2015learning}, AggrevaTeD \cite{sun2017deeply}, DART \cite{laskey2017dart}, and THOR \cite{sun2018truncated}. \sail \cite{bhardwaj2017learning} presents an AggreVaTe-based algorithm for learning heuristic search. We extend \sail by incorporating a recurrent $Q$-like function, in which sense our algorithm more closely resembles AggreVaTeD by Sun \etal \cite{sun2017deeply}. While a recurrent policy can be easily incorporated in AggreVaTeD, we cannot use a policy to evaluate actions. This is due to the fact that we would either have to evaluate all actions in a state, which is computationally infeasible, or we would have to give up on taking actions that are not in the most recent version of the search fringe, which would degrade the performance (see Section \ref{sec:approach}). %RNNs have been previously incorporated in SEARNN by Leblond \etal \cite{searnn} in the context of structured prediction, however we focus on imitation of cost-to-goal oracles. 

%nstead, we propose to use backpropagation through time on the roll-out phase of our algorithm, with cost-to-go supervision on each time step.